\section{引言：不要把编译器画成一个黑箱}
\label{sec:introduction}

在命令行输入一次 \texttt{clang program.c}，得到一个可以运行的文件，很容易造成一种过于
粗糙的心智模型：源程序进入一个名为“编译器”的黑箱，机器程序从另一端出来。这个模型
没有错到不能使用，却粗糙到无法回答任何真正的系统问题。例如，宏为何不会成为抽象语法
树中的节点？局部变量何时不再对应一块内存？外部函数的地址在目标文件中尚未确定，调用
指令又如何能够先被编码？优化后的程序表面上已面目全非，我们依据什么称其“仍是同一个
程序”？这些问题分别属于不同的表示、不同的责任边界，也需要不同种类的证据。

本文采用一个更有解释力的观点：编译系统是一条\emph{表示契约}
（representation contract）链。每个阶段接收一种表示，确立或利用其中的事实，把其中一部分
决策显式化，同时把尚不能决定的问题连同足够的信息交给下一阶段。这里的“契约”不是指每个
内部数据结构都是稳定的公共接口；恰恰相反，Clang AST、LLVM 的优化流水线和机器级中间表示
都可能随版本改变。它指的是更基本的责任关系：词法阶段要保存可供语法分析的记号次序，语义
分析要把名字与声明、表达式与类型联系起来，LLVM IR 要明确控制流与数据流，目标文件要以符号
和重定位描述尚未完成的地址决定，而最终的可执行文件必须在约定输入上实现源程序所要求的可
观察行为。

这种看法也避免了另一种常见误解：教科书中的方框未必对应操作系统中的进程。Clang 的
\emph{驱动程序}（compiler driver）先构造\emph{动作图}（action graph），再把动作绑定为实际
工具与\emph{作业}（job）；预处理、编译后端和集成汇编器可以在同一个
\texttt{clang -cc1} 作业内完成。因而“预处理器、编译器、汇编器、链接器”仍是有用的职责划分，
却不能被误读为“四个必然独立运行的程序”\cite{front-clang-driver,front-clang-command}。本文同时
保留概念分层和实现现实：前者帮助推理，后者防止把一幅示意图冒充运行轨迹。

\subsection{研究对象、问题与边界}

研究对象是 Clang/LLVM 编译系统及其面向 RV64 Linux 的代码生成路径，语言侧以课程规定的
SysY 2022 子集为规范语义载体。全文只追踪一个案例：两个长度均为 8 的整型数组；程序通过
\texttt{getint} 读入前缀长度 $n$，将其限制到 $[0,8]$，逐项相乘后把每个乘积截断到
$[-8,12]$，求和并以 \texttt{putint}、\texttt{putch} 输出。这个例子小到可以逐节点、逐指令
阅读，又同时含有数组、下标、赋值、条件、循环、用户函数与外部运行时调用。它不是为了展示
“语法功能越多越好”，而是为了让同一个事实在多层表示中反复出现。

围绕这一个计算，本文回答三个问题：
\begin{enumerate}
  \item 源字符经过预处理、词法、语法和语义分析后，哪些原本隐含的事实成为带类型的结构？
  \item 结构化控制和可变状态如何依次变成 LLVM IR 的控制/数据流、RV64 指令、ELF 符号与
        重定位，最后成为进程的可观察行为？
  \item “同一程序”的判断可以得到多强的证据？编译成功、结构检查、已知答案测试与形式化证明
        分别能说明什么，又不能说明什么？
\end{enumerate}

本文的范围有意收紧。我们不把若干互不相关的小程序拼成语言特性清单，不把优化后文件变小当作
性能结论，也不声称库存 Clang 是一个符合 SysY 全部规范的前端。纯 SysY 文件是规范性
（normative）源程序；一个只增加 C 声明和预处理指令的 C 观察载体，使 Clang 能展示预处理、
记号与 AST。手写 LLVM IR 和手写 RV64 汇编是对同一语义函数的独立实现，而 Clang 生成的
IR/汇编只作为观察编译器决策的证据。这些身份若混在一起，后续所谓“等价”便没有清楚的比较
对象。

\subsection{方法：沿着一个事实走，而不是沿着工具名走}

本文在每个边界重复询问五件事：输入表示是什么；本阶段确立什么事实；输出表示是什么；案例中
哪一个具体对象承载该事实；还有什么留待下一阶段决定。表~\ref{tab:evidence-kinds} 给出全文使用
的证据标记。其目的不是给叙述增加形式感，而是避免把文档推断写成实验结果，或者把六个测试
样例写成普遍正确性证明。

\begin{table}[t]
\caption{本文的证据等级与允许的结论。}
\label{tab:evidence-kinds}
\centering
\small
\begin{tabular}{p{1.8cm}p{4.2cm}p{5.5cm}}
\hline
标记 & 来源 & 可以支持的结论 \\
\hline
观察 & 固定版本、命令与输入产生的记号、AST、IR、目标文件或运行结果 & 该环境下确实出现了所述结构或行为；不能推出所有版本都相同 \\
文档 & 语言规范、ABI、LLVM/Clang 官方手册 & 接口语义或设计责任；滚动文档仍需与安装版本对应 \\
推断 & 由观察和文档共同给出的因果解释 & 在明确假设下最合理的解释；不是工具实现对未来版本的保证 \\
验证 & 汇编/IR 验证器、结构检查和同一测试预言机 & 排除被检查到的错误；有限测试不等于全域等价证明 \\
\hline
\end{tabular}
\end{table}

这套方法刻意接近《深入理解计算机系统》的写法：不把“编译”当作一个名词，而是指着一个
对象问它是什么。第一个数组元素 \texttt{-4} 在词法层是一元负号与整数常量，在 AST 中是带
类型的一元表达式，在 LLVM 中成为 \texttt{i32 -4} 的初始化数据，在目标文件中成为四个按
目标字节序排列的字节；\texttt{getint} 从已解析的函数调用变成外部声明、未定义符号和调用重
定位，直至链接器从 SysY 运行时库中选择定义。单个对象跨层的连续身份，比罗列每个工具的全部
选项更能暴露系统设计。

\subsection{贡献与文章结构}

本文是一篇实例锚定的技术综述，不提出新的编译算法。它的贡献在于：（1）给出一个无有符号
溢出和越界歧义的 SysY 整数案例及数学契约；（2）把预处理、前端、LLVM 中端、RISC-V 后端、
汇编、ELF 链接和运行时置于同一条可检查的证据链；（3）严格区分驱动动作、实际作业与编译器
内部阶段；（4）以行为预言机和结构证据校准“等价”的强度；（5）在展望中把单层 LLVM IR 与
MLIR/AscendNPU IR 的渐进式降低联系起来，而不把未运行的硬件路径写成实验结果。

第~\ref{sec:case-study} 节固定程序的语义和旅程地图；第~\ref{sec:frontend} 节从字符走到带类型
AST；后续章节依次追踪 LLVM IR、优化与 RV64 代码生成、汇编和链接、端到端验证，并讨论 MLIR
及 AscendNPU IR。全文最终回到同一原则：抽象层不是为了隐藏事实，而是为了让恰当的事实在恰当
的阶段成为可操作对象。
